\chapter{Identificación de requisitos}
\label{cap:requisitos}

Los requisitos se derivan del proceso observado en Innokey, de los objetivos del TFM y de la necesidad de mantener la inteligencia artificial dentro del flujo operativo. Cada requisito incluye un criterio verificable para evitar descripciones demasiado abiertas.

\section{Requisitos funcionales}

\begin{description}
    \item[\textbf{RF1. Carga de la captura.}] La herramienta debe admitir una fotografía estática de la zona de trabajo tomada con una cámara de consumo. Se verifica cuando la aplicación abre la imagen, corrige su orientación y la muestra al usuario.

    \item[\textbf{RF2. Detección de la hoja.}] El sistema debe localizar los cuatro marcadores de esquina y ordenar sus posiciones. Se verifica cuando identifica las cuatro referencias válidas o informa de forma explícita que la rectificación no puede continuar.

    \item[\textbf{RF3. Rectificación geométrica.}] La captura debe transformarse a un plano A4 mediante la homografía calculada con los marcadores. El lienzo resultante debe medir $2100 \times 2970$ píxeles, a una escala de $10\,\text{px/mm}$.

    \item[\textbf{RF4. Detección del WCS.}] La aplicación debe localizar la marca en L grabada por el láser y recuperar su origen y direcciones de eje. Cuando la marca no exista o sea insuficiente, el sistema debe continuar en modo de análisis y bloquear la exportación en coordenadas físicas.

    \item[\textbf{RF5. Localización aproximada de instancias.}] El módulo clásico debe proponer cajas para los toppers mediante componentes conectadas y filtros de área. Estas cajas se utilizan como prompts y no como contorno final.

    \item[\textbf{RF6. Segmentación mediante IA.}] SAM 2 debe producir la máscara binaria final de las instancias localizadas. Si el checkpoint, las dependencias o los prompts no están disponibles, la aplicación debe mostrar un error y no sustituir silenciosamente la salida por la máscara clásica.

    \item[\textbf{RF7. Vectorización.}] Los contornos obtenidos de la máscara de SAM 2 deben simplificarse y transformarse a unidades físicas, conservando una trayectoria exterior por topper. Los cortes interiores solo se habilitan de forma explícita para piezas que realmente los requieran \cite{DouglasPeucker1973}.

    \item[\textbf{RF8. Exportación DXF.}] El sistema debe generar un archivo DXF estructuralmente verificable cuando exista un WCS válido. Las trayectorias de corte y la referencia auxiliar deben almacenarse en capas diferenciadas para conservar la vía de incorporación al flujo CAD/CAM.

    \item[\textbf{RF9. Revisión visual.}] La interfaz debe mostrar la imagen original, la hoja rectificada, la salida derivada de la máscara de SAM 2 y la previsualización vectorial. Las cajas utilizadas como prompts deben conservarse como artefacto de diagnóstico y su procedencia debe quedar registrada en el informe de ejecución.

    \item[\textbf{RF10. Informe de ejecución.}] Cada procesamiento completado debe registrar el método de segmentación, la fuente de los prompts, el estado del WCS, el número de contornos y el tiempo de ejecución.
\end{description}

\section{Requisitos no funcionales}

\begin{description}
    \item[\textbf{RNF1. Bajo coste de captura.}] La solución debe funcionar con un teléfono inteligente y un soporte removible, sin exigir una cámara industrial fija.

    \item[\textbf{RNF2. Reproducibilidad.}] Las particiones de datos, configuraciones, versiones de dependencias y checksums de los modelos deben quedar documentados.

    \item[\textbf{RNF3. Trazabilidad algorítmica.}] La máscara que recibe el extractor de contornos debe ser exactamente la generada por SAM 2. Esta relación debe poder verificarse mediante pruebas automatizadas.

    \item[\textbf{RNF4. Robustez controlada.}] El sistema debe conservar las ocho instancias en las capturas reales evaluadas bajo sombras, recolocación de cámara y cambios de posición del material. Las condiciones que excedan este alcance deben declararse como limitaciones.

    \item[\textbf{RNF5. Interoperabilidad.}] La salida debe utilizar el formato DXF y capas diferenciadas para conservar la compatibilidad estructural con el flujo CAD/CAM. La importación en el programa CAM de Innokey y el corte físico quedan fuera de la validación realizada.

    \item[\textbf{RNF6. Usabilidad.}] El operario debe poder ejecutar el flujo principal desde una sola ventana y recibir mensajes claros cuando falte el WCS o no pueda cargarse el modelo.

    \item[\textbf{RNF7. Ejecución local.}] La segmentación debe poder realizarse en CPU y sin enviar imágenes del taller a servicios externos.

    \item[\textbf{RNF8. Mantenibilidad.}] La detección geométrica, la segmentación, la extracción de contornos y la exportación deben permanecer en módulos separados, de modo que un componente pueda sustituirse sin reescribir el resto del flujo.
\end{description}

\section{Criterios de aceptación del prototipo}

Los criterios de aceptación se expresan como contratos de entrada, procesamiento y salida. Una captura real con cuatro marcadores debe rectificarse y dar lugar a ocho cajas y ocho componentes obtenidas con SAM 2. En el interior del flujo ha de quedar trazado que la máscara neuronal llega al extractor sin sustituciones. Si el WCS es válido, el contrato de salida exige un DXF que pueda reabrirse y contenga ocho polilíneas exteriores cerradas en las capas previstas, con coordenadas finitas, área no nula y límites físicamente plausibles.

En conjunto, estas comprobaciones recorren el software de extremo a extremo y confirman la participación de la inteligencia artificial en la ruta operativa. En lo relativo al DXF, se limitan a la coherencia estructural interna del archivo. El corte físico y la importación en el programa CAM de Innokey requieren una validación posterior.
